\section{A Reporting-and-Decision Protocol}
\label{sec:protocol}

Reporting prescriptions exist (TEE's checklist \citep{messing2026tee}, NIST AI
800-3's repeat-trial minimum \citep{nist2026ai8003}), but not the decision
rule that turns reported uncertainty into a licensed claim. The protocol below
adds that step and nothing more.

When reporting an improvement measured by an LLM judge:
\begin{enumerate}
  \item \textbf{Repeat}: score with $N \geq 3$ independent repeats per
    system in the stated environment.
  \item \textbf{Interval}: report a confidence interval for the gain
    $\Delta$, not a point estimate alone.
  \item \textbf{Decide}: claim only when $\Delta > \mdi_{\alpha}(\env, N)$
    \emph{on the stated item set} \textbf{and} $\fip(\Delta, \env, N) \leq
    \alpha$. The first is a pre-experimental threshold, the second prices the
    gain actually observed; a $\Delta$ just past $\mdi$ can still reverse a
    quarter of the time (Table~\ref{tab:fip-conversion}).
\end{enumerate}
Steps 2 and 3 are both necessary: a gain clearing $\mdi$ whose interval still
covers zero is not licensed.
The two enter at different times: $\mdi$ before the run, to choose the $N$
a target gain requires, and $\fip$ after it, to price the gain observed.
Since the $\fip$ read-off exceeds $\mdi$ in every cell of our grid
(Section~\ref{sec:exp1}), it is usually the $\fip$ condition that binds.
$\fip$ costs no second evaluation: it is estimated from the same repeated
scoring that yields $\mdi$, with repeat~$0$ held out as the screening pass
(Experiment~2), and a team reads its own environment's row off the resulting
conversion table (Table~\ref{tab:fip-conversion}) instead of recomputing one
per decision. None of the three speaks to where the
instrument is aimed (Section~\ref{sec:limitations}). The protocol fixes the \emph{rule}, not a
universal threshold; no constant survives even our own grid, and the
hand-picked ones the loop of Section~\ref{sec:exp4} shipped (a $0.0$ default,
a $+0.02$ remedy) sit below its measured band. %
Measuring your own costs under
\$0.20 per system (dense) or \$10.38 (anchor) at $N{=}20$ over $100$ items;
an $A$-vs-$B$ pair is twice that (Table~\ref{tab:cost}).

\paragraph{Absolute decisions: the SLA margin table.}
The rule above is relative ($A$ vs.\ $B$); delivery decisions are often
absolute, and the same null distribution prices a contract threshold $t$ too
(Table~\ref{tab:sla-margin}): G-theory's absolute/relative distinction
\citep{cronbach1972dependability, messing2026tee}.
Every entry is \emph{fixed-prompt}, and Section~\ref{sec:exp3}'s shared prompt
component exceeds every margin in the table
(Appendix~\ref{app:extra}): re-wording a rubric moves an absolute score by more
than the margin it must clear. Re-calibrate rather than reuse.



Every number in this paper passes all three steps, traceable to a derived
record and a run identifier.
